% Lösungen Punkte und Strecken im Koordinatensystem · Streckenlänge, Mittelpunkt, Teilpunkte · Prüfungs-Fokus Abitur GK (zusammenbau v0.6)
\einheitenkopf{Einheit 2 · Streckenlängen, Mittelpunkt und Teilpunkte}

\erg{1}{a) ein Punkt – die Mitte von $AB$ \quad b) $\overrightarrow{AB} = (2 | 3 | 6)$, $|\overrightarrow{AB}| = \sqrt{4 + 9 + 36} = 7$ LE \quad c) $|\overrightarrow{SH}| = \sqrt{4 + 9 + 36} = 7$ m; mit Zuschlag $0{,}4$ m (nicht $40$ m): $4 \cdot 7{,}4 = 29{,}6$ m}
\erg{2}{a) $|\overrightarrow{AB}| = \sqrt{36 + 64 + 0} = 10$ m; vorher $10 : 1{,}25 = 8$ m (nicht $10$ minus $25\,\%$) \quad b) $|\overrightarrow{AB}| = \sqrt{4 + 36 + 81} = 11$ m; vorher $11 : 1{,}1 = 10$ m (nicht $11$ minus $10\,\%$) \quad c) $|\overrightarrow{AB}| = \sqrt{1 + 16 + 64} = 9$ cm; vorher $9 : 1{,}5 = 6$ cm (nicht $9$ minus $50\,\%$) \quad d) $\frac12 \cdot (\overrightarrow{OD} + \overrightarrow{OS}) = (2 | 1{,}5 | 3)$, also ist $W$ Mittelpunkt; an der $x_1x_3$-Ebene gespiegelt: $T(2 | -1{,}5 | 3)$ \quad e) $\frac12 \cdot (\overrightarrow{OD} + \overrightarrow{OS}) = (3 | 3 | 4)$, also ist $W$ Mittelpunkt; an der $x_2x_3$-Ebene gespiegelt: $T(-3 | 3 | 4)$ \quad f) $\frac12 \cdot (\overrightarrow{OD} + \overrightarrow{OS}) = (4 | 3 | 2)$, also ist $W$ Mittelpunkt; an der $x_1x_2$-Ebene gespiegelt: $T(4 | 3 | -2)$ \quad g) $T_1 = P + \frac{1}{3} \cdot \overrightarrow{PQ} = (6 | 1 | 2)$, $T_2 = P + \frac{2}{3} \cdot \overrightarrow{PQ} = (3 | 2 | 1)$ (nicht die Mitte) \quad h) Anteil $\frac{2}{5}$ (nicht $\frac{2}{3}$): $G = A + \frac{2}{5} \cdot \overrightarrow{AC} = (5 | 2 | 4)$ \quad i) Anteil $\frac{3}{4}$ (nicht $\frac{1}{4}$ – $G$ liegt näher an $C$): $G = A + \frac{3}{4} \cdot \overrightarrow{AC} = (5 | 6 | 2)$ \quad j) $D = A - \overrightarrow{AB} = (1 | 4 | 0)$ (nicht $B$ und nicht $A + 2\overrightarrow{AB}$) \quad k) $D = A - \overrightarrow{AB} = (-2 | -6 | 4)$ (nicht $B$ und nicht $A + 2\overrightarrow{AB}$) \quad l) $D = A - \overrightarrow{AB} = (-2 | 5 | 4)$ (nicht $B$ und nicht $A + 2\overrightarrow{AB}$)}
\erg{3}{a) $\overrightarrow{AC} = \overrightarrow{OC} - \overrightarrow{OA} = 2 \cdot \overrightarrow{OA} = (4 | -2 | 4)$, $|\overrightarrow{AC}| = 6$ (nicht $|\overrightarrow{OC}|$) \quad b) $\overrightarrow{AC} = \overrightarrow{OC} - \overrightarrow{OA} = -2 \cdot \overrightarrow{OA} = (-4 | -6 | -12)$, $|\overrightarrow{AC}| = 14$ (nicht $|\overrightarrow{OC}|$) \quad c) $\overrightarrow{AC} = \overrightarrow{OC} - \overrightarrow{OA} = -\frac{1}{2} \cdot \overrightarrow{OA} = (-2 | -2 | -1)$, $|\overrightarrow{AC}| = 3$ (nicht $|\overrightarrow{OC}|$) \quad d) Tunnellänge $|\overrightarrow{AB}| \cdot 10$ m $= 500$ m; gleiche Fahrzeit: $\frac{s}{90} = \frac{500 - s}{60}$ ergibt $s = 300$ m (Weg des Zugs ab $B$). \quad e) $|\overrightarrow{PQ}| = 10$, also $1000$ m; $\frac{s}{15} = \frac{1000 - s}{10}$ ergibt $s = 600$ m; Treffpunkt $P + 0{,}6 \cdot \overrightarrow{PQ} = (5{,}6 | 5{,}8 | 0)$}
